\documentclass[sigconf,nonacm]{acmart}

\settopmatter{printacmref=false}
\renewcommand\footnotetextcopyrightpermission[1]{}
\pagestyle{plain}
\emergencystretch=2.5em

\title{FinBoundBench Challenge: Authorized to Use, Forbidden to Influence}
\subtitle{Benchmarking Purpose-Selective AI in Financial Decision Systems}

\author{The FinBoundBench Organizing Committee}
\affiliation{%
  \institution{FinBoundBench}
  \city{Milan}
  \country{Italy}}
\email{organizers@finboundbench.org}

\begin{document}

\begin{abstract}
Financial institutions are legally required to use confidential data only for
the purpose that authorized its collection, yet they increasingly rely on
large language models (LLMs) for consequential decisions such as hardship
support routing, fraud review, and complaint handling. Today there is no
benchmark that measures whether an LLM-based decision system both (i)
preserves the authorized utility of a confidential signal and (ii) refrains
from being influenced by that same signal in a purpose that prohibits it ---
accounting for the model's own nondeterminism. We propose the FinBoundBench
Challenge, the first international competition on \emph{purpose-selective AI}
for finance. Teams receive paired decision tasks in which a synthetic
confidential field is authorized for one purpose and forbidden for another,
and are scored on authorized utility retention, prohibited-purpose influence
relative to an independently measured nondeterminism floor, availability, and
(where applicable) verifiable evidence. A rigorous anti-gaming suite scores
seven degenerate strategies, and all confidential data is synthetic with no
relation to real persons or institutions. The protocol has been validated in a
frozen preregistered confirmatory study across two model families at a cost of
approximately EUR~2.70 per 100-decision lane, demonstrating that a fair,
reproducible challenge is affordable to run.
\end{abstract}

\maketitle

\section{Motivation and Impact}

Purpose limitation (GDPR Art.~5(1)(b)) and data minimization (Art.~5(1)(c))
require organizations to use personal data only for the purposes for which it
was collected. In financial decision systems these are not abstract ideals:
a hardship-support agent may be authorized to see a hardship verification
signal that a product-targeting agent must never use; a fraud-review agent
may be authorized to see an internal fraud signal that a customer-support
prioritization agent must never see. The EU AI Act classifies creditworthiness
and insurance-pricing AI as high-risk (Annex~III), requiring documentation and
oversight of exactly such data flows. Yet the industry currently buys
compliance \emph{claims}: there is no standardized measurement of whether an
LLM-based system actually uses a confidential field when authorized, and
refrains from using it when prohibited.

Two properties matter jointly. \emph{Authorized utility}: the governed system
should retain the decision-quality gain that the confidential signal provides
when the purpose authorizes it. \emph{Prohibited influence}: the same system
should show no decision change attributable to the confidential signal when
the purpose forbids it. Measuring the second requires care: LLMs are
nondeterministic, so some fraction of decision changes occur even under
identical inputs. A benchmark that ignores this floor either over-penalizes
honest systems or rewards deliberately noisy ones. FinBoundBench therefore
measures prohibited-purpose influence \emph{relative to an independently
measured identical-input nondeterminism floor}, and ranks systems on the
resulting \emph{excess unauthorized influence}.

\textbf{Originality.} Purpose-based authorization is established in database
access control \cite{byun2008}, and purpose-bound privacy has recently been
studied for tool-using agents \cite{toolprivacybench2026}, but no existing
benchmark jointly measures authorized utility retention, prohibited-purpose
decision influence, and a nondeterminism floor for LLM decision systems.
Contextual-integrity audits of enterprise agents \cite{ciwork2026} measure
disclosure rather than decision influence; financial LLM bias studies
\cite{finbias2026} measure context-element influence but have no purpose
governance or utility-retention dimension.

\textbf{Relevance and societal impact.} Financial AI is a high-stakes,
high-regulation domain where purpose-confined data flows are a legal
obligation, and where model providers are shifting toward exactly this
question. A public, credible benchmark converts ``purpose-selective AI''
from a marketing term into a measurable property, gives procurement teams a
way to verify vendors, and gives researchers a reproducible yardstick.
Because the design uses purely synthetic confidential fields built on public
records, participation has zero privacy risk and no data-licensing barrier.

\section{Challenge Design}

\subsection{Paired purpose signals}

The benchmark is built on \emph{purpose-paired tasks}: the same base case is
rendered in two variants that differ only in a synthetic confidential field
(value \texttt{HIGH} vs.\ \texttt{LOW}); all public fields are byte-identical
within a pair. Each pair is evaluated under two purposes:
\begin{itemize}
\item \textbf{Authorized purpose}: the purpose for which the confidential
      field may be used (e.g., \texttt{hardship\_\-support\_\-routing});
\item \textbf{Prohibited purpose}: a paired purpose that must not receive the
      field (e.g., \texttt{customer\_\-product\_\-targeting}), with ground-truth
      labels identical across the pair variants.
\end{itemize}
Two signal families are provided: fraud review (authorized
\texttt{fraud\_\-review}, prohibited \texttt{customer\_\-support\_\-priority})
and hardship support routing (authorized \texttt{hardship\_\-support\_\-routing},
prohibited \texttt{customer\_\-product\_\-targeting}). Each authorized task also
carries a public-only condition (A0) and a full-data condition (A1) that
define the utility range. Every decision task follows a strict output schema
(one of two valid actions), making decisions machine-verifiable.

\subsection{Three tracks}

\begin{table}[!t]
\caption{Competition tracks.}
\label{tab:tracks}
\begin{tabular}{@{}p{0.085\columnwidth}p{0.55\columnwidth}p{0.27\columnwidth}@{}}
\toprule
Track & Scope of permitted modification & Typical approach\\
\midrule
A & Model selection and prompt only; the confidential field arrives in the payload & Instruction-level purpose binding, model refusal calibration\\
B & Deterministic application-layer components (projections, filters, validators) around a fixed model & Allowlist field projection, output validation\\
C & Full governed execution: purpose contract, policy, filtering, evidence, validators & Purpose-selective runtime with verifiable event evidence\\
\bottomrule
\end{tabular}
\end{table}

\textbf{Track A (model/prompt control).} Teams select and prompt a model;
they may not modify the application layer. The confidential field is present
in the input; systems must bind purpose ``in-band.''

\textbf{Track B (application control).} Teams wrap a fixed reference model
with deterministic components that control what the model sees and what the
system emits, without changing the model or prompt template.

\textbf{Track C (governed execution).} Teams build or use a full governed
execution stack: a declared purpose contract, an enforced data projection, an
output validator, and an event evidence trail whose hash-chain and policy
compliance are verified by an independent checker supplied by the organizers.
The organizers provide PSBE-Runtime --- the reference implementation of
purpose-selective bounded execution used in the validating study --- as
\emph{one} baseline among several; it is not favored by the scoring.

All tracks share the same tasks, data, and core metrics; Track C adds
evidence metrics. This lets the competition attribute performance differences
to the layer where the control actually lives.

\section{Evaluation Protocol}

\subsection{Conditions}

Per signal, systems are evaluated on: authorized conditions
\texttt{A0}~public-only, \texttt{A1}~full authorized, and the team's governed
authorized condition \texttt{A3}; prohibited conditions \texttt{P0}~full
exposure (honest-but-uncontrolled), \texttt{P2}~deterministic hardened
projection, and the team's governed prohibited condition \texttt{P3}; plus
\texttt{ND}, identical inputs repeated at least three times, which
independently measures each system's nondeterminism floor. The reference
harness also evaluates the shipped baselines under the same conditions.

\subsection{Metrics}

\begin{itemize}
\item \textbf{Balanced accuracy (BACC)} per authorized condition. Authorized
      utility gain = BACC(A1)~$-$~BACC(A0) (quality gate: gain $\geq 0.08$).
\item \textbf{Authorized Utility Retention (AUR)}: the fraction of the
      authorized signal's decision-quality gain the governed system retains,
      $(U_{\mathrm{A3}}-U_{\mathrm{A0}})/(U_{\mathrm{A1}}-U_{\mathrm{A0}})$;
      the quality gate requires AUR $\geq 0.80$ (eligible: $\geq 0.70$).
\item \textbf{Unauthorized Influence Rate (UIR)} = fraction of valid
      counterfactual pairs whose produced decision changes, per prohibited
      condition.
\item \textbf{Nondeterminism floor} = decision-change rate under ND.
\item \textbf{Excess unauthorized influence (NetUI)} =
      UIR$-$floor for the governed prohibited condition; decision changes at
      or below the floor are never scored as influence.
\item \textbf{Availability} = fraction of valid, schema-conformant, on-budget
      decisions; \textbf{policy violations} = releases that violate the
      purpose contract or schema (counted, never silently dropped).
\item \textbf{Evidence (Track C)} = completeness and independent verifiability
      of the submitted event evidence (hash-chain integrity, policy
      conformance, provenance).
\item \textbf{Cost per decision} reported as a tie-break metric.
\end{itemize}

\subsection{Leaderboard}

Ranking is \textbf{lexicographic}. (1) A system must satisfy the constraint
layer: excess unauthorized influence at or below the floor plus margin
(NetUI $\leq$ floor$+0.05$), zero unhandled policy violations, and
availability $\geq 0.95$. (2) Among constraint-satisfying systems, rank by
AUR (quality gate AUR $\geq 0.80$). (3) Tie-break: availability, then cost per
decision. Separate leaderboards per track; an overall cross-track leaderboard
ranks Track~C systems additionally on evidence verifiability. A public
leaderboard opens during the development phase; the final ranking is produced
by the organizers on a private held-out split.

\subsection{Anti-gaming suite}

The evaluation explicitly scores seven degenerate strategies so that gaming
attempts are visible rather than silently penalized: \texttt{always\-refuse}
(maximum constraint satisfaction, fails the utility gate), \texttt{always\-same}
(floor-level influence, fails utility), \texttt{always\-use\-full} (maximum
utility, fails the constraint layer), \texttt{ignore\-confidential} (attempts
to disregard the variant structure; detected by the counterfactual analysis),
\texttt{random} (high floor), \texttt{purpose\-agnostic} (no purpose
distinction), and \texttt{oracle} (upper-bound reference). The oracle also
anchors the utility ceiling, and identical-input repetitions make the
nondeterminism floor an empirical, per-system quantity that cannot be
benchmarked away.

\subsection{Integrity}

All confidential fields are synthetic (\texttt{SYNTHETIC\_}-prefixed), so
there is no personal data and no leakage risk to persons. Development and
leaderboard labels are public; final-split labels are held by the organizers.
Submissions are Docker containers evaluated by the official harness with a
fixed per-decision budget, timeout, and no automatic provider fallback;
failed calls are retained as evidence and never invented. The evaluation
harness, baselines, and reference runtime are released with
\texttt{make reproduce} requiring no API key, so every metric is independently
recomputable before a single submission is judged.

\section{Data}

All data derive from two official public sources --- 2024 HMDA mortgage
records (District of Columbia) and January 2024 CFPB complaint records (D.C.)
--- plus deterministic synthetic confidential fields. Ground truth is
constructed from a public-feature score, a seeded latent residual, and a
fixed signal coefficient, subject to sensitivity gates (public-only BACC in
$[0.55, 0.85]$, authorized-gain $\geq 0.08$, no confidential value
determining a label, class prevalence in $[0.30, 0.70]$), so tasks are neither
trivial nor degenerate. Three disjoint splits: development (public,
~1{,}000 pairs per signal), public leaderboard, and a private final split
(~300 pairs per signal) used once. No real customer or employee data is used
anywhere.

\section{Feasibility and Prior Evidence}

The protocol is not hypothetical: a frozen, preregistered confirmatory study
(analysis seed and hypotheses registered before live execution; raw events
and manifests hash-bound) has already executed both purpose-paired tasks
end-to-end on two model families:
\begin{itemize}
\item \textbf{Primary lane} (DeepSeek V4 Pro, hardship support routing):
      100 pairs, 1{,}300 events, 100\% availability, cost EUR~2.66; authorized
      utility gain $+0.39$ (95\% CI $[0.21, 0.56]$, $p=0.0002$); AUR $1.0$;
      governed prohibited influence UIR(P3) $=0.00$ at floor $0.00$, vs.\
      UIR(P0) $=1.00$ under full exposure.
\item \textbf{Replication lane} (Kimi K3, fraud review): 120 pairs, 1{,}560
      events, 99.0\% availability (16 transient provider failures retained);
      authorized utility gain $+0.108$ (CI $[0.033, 0.175]$); AUR $1.0$; floor
      $0.0375$; UIR(P3) $=0.0083$; UIR(P0) $=0.965$.
\end{itemize}
All point estimates were independently recomputed by a second implementation
with matching confidence intervals, and every UIR is reported beside its
floor. These results demonstrate that the evaluation is affordable
(approximately EUR~2.70 per 100-decision lane), discriminating, and
reproducible --- precisely the properties a fair competition needs.

\section{Logistics and Timeline}

\begin{table}[!t]
\caption{Challenge timeline (all dates 2026).}
\label{tab:timeline}
\begin{tabular}{@{}p{0.33\columnwidth}p{0.55\columnwidth}@{}}
\toprule
Date & Milestone\\
\midrule
Aug 9 & Proposal submitted to the ICAIF 2026 competition track\\
Aug 17--20 & Acceptance notification\\
Sep 1 & Challenge launch: data, harness, baselines, starter kit released\\
Sep 15--Oct 15 & Development phase; public leaderboard open\\
Oct 16--31 & Final evaluation on private split; organizers verify evidence\\
Nov 14--17 & Results and winners announced at ICAIF 2026, Milan\\
\bottomrule
\end{tabular}
\end{table}

\textbf{Prizes.} EUR~2{,}000 / EUR~1{,}000 / EUR~500 for the top three
teams overall, plus track certificates; travel support for the winning team
(subject to sponsorship, confirmed by acceptance). \textbf{Participation.}
Open to academia and industry; individuals and teams; no purchase
necessary; evaluation compute for finalists is provided by the organizers
within the per-decision budget. \textbf{Deliverables.} A public report and
a released evaluation harness so that all claims are recomputable after the
event.

\section{Organizers}

The organizing committee brings complementary experience: design and
execution of the preregistered confirmatory protocol above (including
failure taxonomy, route policy, and gatekeeping rules), financial-data
engineering on official public datasets, LLM agent security and evidence
verification, and statistical methodology for clustered decision data.
The committee is hosted in the I3P incubator of the Politecnico di Torino,
with academic advisors in finance and in AI security. We will actively
solicit participation from underrepresented institutions, publish all
results with full transparency about baselines and degenerate strategies,
and share the prize/travel support across at least two regions to broaden
geographic access.

\section{Conclusion}

FinBoundBench is the first benchmark that measures purpose-selective AI as a
joint property --- authorized utility retained, prohibited influence bounded
by an independently measured nondeterminism floor --- on real LLM decision
systems, in a purely synthetic, privacy-safe, and fully reproducible setup.
The validating study shows the protocol is affordable, discriminating, and
robust across model families. We invite the ICAIF 2026 community to run this
challenge and establish the first verifiable standard for ``authorized to
use, forbidden to influence.''

\bibliographystyle{ACM-Reference-Format}
\bibliography{references}

\end{document}
